\chapter*{Preface}
The integers are the oldest object of mathematics and still the most mysterious. A child can state their open
problems; the proofs of their theorems call on the whole of mathematics. This volume of \emph{Order and Chance}
studies them with elementary tools, and finds at every turn that elementary does not mean simple.

It starts with divisibility and the Euclidean algorithm, the primes and their infinitude, and congruences, the
arithmetic of remainders. Arithmetic functions and their Dirichlet convolution give the first glimpse of the
analytic theory of Tome~5. The multiplicative group modulo a prime is cyclic, and the squares in it obey a law
that Gauss called the \emph{theorema aureum}: quadratic reciprocity, the centrepiece of the book, proved here more
than once because each proof reveals a different reason. From there the book follows questions that look
innocent: which integers are sums of two or four squares, which equations $x^2-dy^2=1$ have solutions and how
continued fractions find them all, which Diophantine equations can be solved at all. To answer them it enlarges the
integers, to the Gaussian and Eisenstein integers, to ideals, where unique factorisation is restored after it
fails, and finally to the $p$-adic numbers, a completion of $\Q$ in which divisibility by $p$ becomes nearness.

The goal is to see the integers as a structure: their arithmetic is governed by groups, rings and, already, by
analysis. The same arithmetic runs modern cryptography, and the companion Python notebooks follow it there, from
fast modular exponentiation to RSA and primality tests.

Every theorem, proposition and lemma is proved; results quoted only for their interest appear as remarks, with a
reference to where they are proved. Exercises are graded \textbf{$\star$} for a warm-up, \textbf{$\star\star$} for a good
afternoon and \textbf{$\star\star\star$} for gems, and every one has a hint and a complete solution at the end of its
chapter, beginning with \emph{How to find it} before the \emph{Write-up}. ``Pause and think'' boxes ask you to guess
before you read on.
